\chapter{État de l'art}

\section{Introduction du chapitre}

La reconnaissance automatique des langues des signes consiste à analyser des productions visuelles afin d'en identifier les unités linguistiques. Ce chapitre présente les propriétés de ces langues, les principales approches de capture et de représentation, puis les familles de modèles et des travaux représentatifs. Il se termine par une synthèse des tendances observées.

\section{La LSM comme langue multi-articulateurs}

Comme les autres langues des signes, la Langue des Signes Malgache (LSM) mobilise plusieurs articulateurs visuels et spatiaux, dont les mains, le visage, le regard et le haut du corps. Ces éléments peuvent transmettre simultanément des informations linguistiques, contrairement à la parole, qui se déploie principalement dans le temps par le canal vocal-auditif \cite{Sandler2006SignLanguageUniversals}.

Les composants manuels comprennent la configuration et l'orientation des mains, leur emplacement ainsi que leur mouvement. Lorsqu'elles sont utilisées ensemble, les deux mains peuvent jouer des rôles distincts : l'une peut notamment établir un repère spatial que l'autre réutilise. La fonction dépend du signe et du contexte ; elle n'est pas identique dans toutes les productions.

Les composants non manuels comprennent notamment les expressions du visage, le regard et les mouvements de la tête ou du tronc. Ils accompagnent la production manuelle et peuvent contribuer à l'organisation grammaticale, à la prosodie et à l'expression des intentions du signeur. La description d'une langue des signes doit donc considérer la coordination de ces composantes, plutôt que réduire le signe à une forme de la main isolée.

\section{Approches de capture et d'extraction}

La capture des signes peut s'appuyer sur des capteurs portés, tels que des gants instrumentés ou des capteurs inertiels, et sur des dispositifs optiques comme les caméras RGB ou de profondeur. Les capteurs portés mesurent des mouvements ou des configurations au moyen d'instruments attachés au corps ; les caméras observent la production sans équipement porté. Ces modalités impliquent des compromis différents en matière de matériel, de liberté de mouvement et de type d'information recueillie.

À partir d'une vidéo, les méthodes peuvent analyser directement les pixels ou en extraire des représentations intermédiaires. Les points de repère décrivent les articulations ou les parties du corps par des coordonnées, tandis que les images brutes conservent aussi l'apparence et le contexte visuel. Les représentations squelettiques sont étudiées pour la reconnaissance de langues des signes et peuvent être associées à plusieurs familles de modèles \cite{Gonfa2025Ethiopian}. MediaPipe est l'un des cadres permettant l'extraction de points de repère du corps et des mains \cite{Lugaresi2019MediaPipe,Zhang2020MediaPipeHands,Bazarevsky2020BlazePose}.

\begin{figure}[ht]
\centering
\begin{minipage}[t]{0.46\textwidth}
\centering
\includegraphics[width=\linewidth]{\detokenize{../figures/images du memoire/poselandmark-mediapipe.png}}
\par\small Points de pose du corps
\end{minipage}
\hfill
\begin{minipage}[t]{0.46\textwidth}
\centering
\includegraphics[width=\linewidth]{\detokenize{../figures/images du memoire/handlandmark-mediapipe.png}}
\par\small Points de repère d'une main
\end{minipage}
\caption{Repères corporels et manuels de MediaPipe.}
\label{fig:mediapipe-landmarks}
\end{figure}

Les coordonnées extraites peuvent ensuite être normalisées afin de limiter l'effet des différences de position, d'échelle ou de morphologie entre signeurs. Les références anatomiques et les transformations retenues varient selon les représentations et les tâches ; elles doivent être précisées et évaluées dans chaque protocole.

\section{Architectures de classification}

Plusieurs familles d'architectures peuvent classifier des séquences de points clés.

\subsection{Réseaux récurrents}

Les LSTM et les GRU sont des réseaux récurrents conçus pour représenter des séquences et leurs dépendances temporelles. Le LSTM introduit une mémoire contrôlée par des portes \cite{Hochreiter1997LSTM} ; le GRU est une autre cellule récurrente proposée dans le cadre des modèles encodeur-décodeur \cite{Cho2014GRU}. Une étude sur la langue des signes indonésienne applique MediaPipe Holistic et compare ces deux familles en temps réel ; elle rapporte 94 \% pour le LSTM et 96 \% pour le GRU \cite{Ilham2024Indonesian}. Ces résultats dépendent du corpus et du protocole propres à cette étude et ne sont pas directement généralisables.

\subsection{Transformers}

Les Transformers utilisent un mécanisme d'attention pour traiter les relations entre éléments d'une séquence \cite{Vaswani2017}. Des travaux ont étudié cette architecture pour la reconnaissance de signes arabes à partir de points de repère \cite{Alyami2023Arabic}. Cette approche constitue une alternative aux modèles récurrents pour l'analyse de séquences.

\subsection{Réseaux denses et algorithmes classiques}

Pour des données statiques ou des vecteurs agrégés, des réseaux denses (MLP) ou des algorithmes classiques comme les SVM peuvent également être étudiés. Ils constituent des points de comparaison possibles lorsque la représentation et le protocole le permettent.

\subsection{Comparaison des architectures}

Comparer plusieurs architectures sur les mêmes données et selon le même protocole permet d'isoler plus clairement l'effet des choix de modélisation. Les résultats restent toutefois dépendants de la taille et de la diversité du corpus, de la représentation choisie et des conditions d'évaluation.

\section{Comparaison de travaux sur les langues des signes}

Les études comparées ci-dessous illustrent plusieurs choix de représentation et de classification pour la reconnaissance de signes. Le tableau retient les caractéristiques utiles pour situer ces approches : langue étudiée, représentation ou modèle, et résultat ou apport méthodologique. Les performances ne sont pas directement comparables, car les corpus et les protocoles diffèrent.

\begin{table}[ht]
\centering
\small
\caption{Études sélectionnées en reconnaissance des langues des signes.}
\label{tab:comparaison-travaux-ls}
\begin{tabular}{|p{3.0cm}|p{4.3cm}|p{6.0cm}|}
\hline
Étude et langue & Approche étudiée & Résultat ou intérêt pour la comparaison \\
\hline
Ilham et al. (2024), langue des signes indonésienne \cite{Ilham2024Indonesian} & MediaPipe Holistic, comparaison LSTM--GRU en temps réel & Les auteurs rapportent 94 \% avec LSTM et 96 \% avec GRU dans leur protocole. \\
\hline
Gonfa et al. (2025), langue des signes éthiopienne \cite{Gonfa2025Ethiopian} & Représentation squelettique ; CNN-LSTM, LSTM, BiLSTM et GRU & Étudie plusieurs architectures sur une même représentation squelettique. \\
\hline
Alyami et al. (2023), langue des signes arabe \cite{Alyami2023Arabic} & Transformer et points de repère pour des signes isolés & Illustre une approche par landmarks associée à une architecture différente des réseaux récurrents. \\
\hline
Maquera et al. (2024), langue des signes péruvienne \cite{Maquera2024LSP} & Reconnaissance avec LSTM et restitution en quechua & Montre un exemple de travaux reliant reconnaissance de signes et langue parlée locale ; aucun score n'est repris ici. \\
\hline
\end{tabular}
\end{table}

Ces travaux montrent que les représentations squelettiques et plusieurs familles de modèles sont étudiées dans des contextes linguistiques distincts. Les résultats rapportés ne peuvent pas être transposés directement d'une langue ou d'un corpus à un autre ; ils servent plutôt à identifier des choix méthodologiques pertinents.

\section{Positionnement du projet}

Les travaux examinés montrent que les résultats dépendent du type de capture, de la représentation des signes, de l'architecture et du protocole expérimental. Dans ce contexte, le mémoire porte sur la comparaison de deux modèles récurrents appliqués à la Langue des Signes Malgache. Les chapitres suivants décrivent les données, le protocole et les limites nécessaires à l'interprétation de cette étude.

\section{Synthèse du chapitre}

Les approches étudiées combinent différentes modalités de capture, représentations visuelles et architectures de classification. Leurs résultats dépendent des données et des protocoles employés, ce qui rend nécessaire une description précise de ces éléments dans toute étude comparative.
